\section{Finite certificates and the dependence on the exponent}\label{rex:sec:certificates}

\subsection{Finite products bracket the exact constant}

The constant in \cref{rex:eq:constants} is expressed by a standard special
function, but it also admits rigorous elementary enclosures.

\begin{proposition}[Product certificates]\label{rex:prop:certificates}
For every $m>0$ and integer $j\geq1$,
\begin{equation}\label{rex:eq:L-brackets}
 \frac jm t_j^{m+1}\leq L_m\leq
 \left(\frac jm+\frac1{m+1}\right)t_j^{m+1}.
\end{equation}
Consequently,
\begin{equation}\label{rex:eq:c-brackets}
 \frac{t_j^{-(m+1)}}{j/m+1/(m+1)}
 \leq c_m\leq\frac mj t_j^{-(m+1)}.
\end{equation}
The ratio of the upper bound to the lower bound for $c_m$ is exactly
$1+m/((m+1)j)$. For positive integer $m$, all endpoints are rational.
\end{proposition}

\begin{proof}
Integrating \cref{rex:eq:f-derivative} between $0$ and $t_j$ gives
\[
 0\leq L_m-f_m(t_j)\leq\int_0^{t_j}t^m\,dt
 =\frac{t_j^{m+1}}{m+1}.
\]
Since $f_m(t_j)=(j/m)t_j^{m+1}$, this proves
\cref{rex:eq:L-brackets}. All quantities are positive, so taking
reciprocals proves \cref{rex:eq:c-brackets}. The endpoint ratio follows
by cancellation. When $m$ is a positive integer, the factors defining
$t_j$ and the powers in these formulas are rational.
\end{proof}

These intervals certify the constant without floating-point Gamma
evaluation. Their relative gap tends to zero at an explicit rate.
They do not by themselves certify how close a finite threshold
$T_m(K)$ is to its asymptotic equivalent: that is a different error
problem.

There is also an exact finite computation bound. If the backward
recurrence has been computed only as far as stage $k$, then
\begin{equation}\label{rex:eq:partial-computation}
 k^m q_k\leq T_m(K)<k^m q_k+\sum_{r=1}^{k-1}r^m
 \quad(k\geq2).
\end{equation}
For integer $m$, this is a certified rational or integer enclosure of a
particular threshold, independent of the asymptotic theorem.

\subsection{An integral interpretation of the constant}

Since $f_m(1)=0$, \cref{rex:eq:f-derivative} also gives
\begin{equation}\label{rex:eq:residual-integral}
 L_m=\int_0^1t^m
 \left[C\!\left(\frac{my_m(t)}t\right)
             -\frac{my_m(t)}t\right]dt.
\end{equation}
The term in brackets is a deterministic rounding residual. Formula
\cref{rex:eq:residual-integral} shows precisely how the successive
increment regimes determine its weighted average. Replacing this
residual by a fixed average would discard the dependence of its phase
on the current quotient.

The full-trajectory theorem gives a related limit for the distribution
of the losses. For a trajectory ending at $J=\tau_m(n)$, define a
probability measure on $[0,1]$ by
\begin{equation}\label{rex:eq:loss-measure}
 \mu_n=\sum_{k=2}^{J}\frac{x_{k-1}-x_k}{n}\,\delta_{k/J}.
\end{equation}
Here $\delta_u$ denotes unit mass at $u$. The weights sum to one because
$x_1=n$ and $x_J=0$.

\begin{corollary}[Limiting distribution of rounding losses]\label{rex:cor:losses}
The measures $\mu_n$ converge weakly to the absolutely continuous
probability measure with distribution function $1-F_m(t)$ and density
\begin{equation}\label{rex:eq:loss-density}
 \rho_m(t)=c_m t^m
 \left[C\!\left(\frac{my_m(t)}t\right)
             -\frac{my_m(t)}t\right]
 \quad\text{for almost every }0<t<1.
\end{equation}
In particular, $0\leq\rho_m(t)\leq c_m t^m$ and
$\int_0^1\rho_m(t)dt=1$. The same conclusion holds under the
smooth-weight hypothesis of \cref{rex:thm:weights-main}.
\end{corollary}

\begin{proof}
The distribution function of \cref{rex:eq:loss-measure} is $1-X_n(t)$.
The uniform convergence of the forward paths therefore proves
convergence of these distribution functions to $1-F_m$. The
derivative formula in \cref{rex:eq:f-derivative}, including its integrable
bound at zero, shows that the limit is absolutely continuous and gives
\cref{rex:eq:loss-density}. For smooth weights, use the forward-profile
part of \cref{rex:thm:weights-main}, proved below.
\end{proof}

No randomness has been introduced in this statement. The measure records
what fraction of a fixed initial value is lost at each normalized stage.

\subsection{Monotonicity and endpoint behavior of \texorpdfstring{$c_m$}{c(m)}}

The explicit constant permits useful parameter calculations. Let
$\gamma$ be Euler's constant and $\zeta$ the Riemann zeta function.

\begin{proposition}[Dependence on $m$]\label{rex:prop:parameter}
The function $m\mapsto c_m$ is strictly increasing on $(0,\infty)$.
Moreover,
\begin{align}
 \log c_m
 &=-m\log m+(1-\gamma)m+O(m^2),
        &&m\downarrow0,\label{rex:eq:small-m}\\
 \log c_m
 &=\log m+\gamma+
   \sum_{r=2}^{\infty}\frac{\zeta(r)}{r(m+1)^{r-1}},
        &&m>0,\label{rex:eq:large-m-series}\\
 c_m
 &=e^\gamma m\left(1+\frac{\pi^2}{12(m+1)}+O(m^{-2})\right),
        &&m\to\infty.\label{rex:eq:large-m}
\end{align}
In particular, $c_m\to1$ as $m\downarrow0$, and $c_m/m$ decreases
strictly to $e^\gamma$ as $m\to\infty$.
\end{proposition}

\begin{proof}
The convergent expansion
\[
 \log\Gamma(1-z)=\gamma z+
      \sum_{r=2}^{\infty}\frac{\zeta(r)}r z^r,
 \qquad |z|<1,
\]
is a standard Gamma identity~\cite{dlmf}. Substitute $z=1/(m+1)$
to obtain \cref{rex:eq:large-m-series}. Its positive terms decrease with
$m$, which proves the assertion about $c_m/m$ and the expansion
\cref{rex:eq:large-m}. For small $m$, put $a=m/(m+1)$ and use
$\log\Gamma(a)=-\log a-\gamma a+O(a^2)$; substitution in
\cref{rex:eq:constants} gives \cref{rex:eq:small-m}.

For strict monotonicity of $c_m$ itself, logarithmic differentiation gives
\[
 \frac{d}{dm}\log c_m
 =\frac1m+\log\Gamma(a)+(1-a)\psi(a),
 \qquad a=\frac m{m+1},
\]
where $\psi$ is the digamma function. Using its functional equation,
this is
\[
 g(a)=\log\Gamma(1+a)-\log a+(1-a)\psi(1+a).
\]
We have $g(1)=0$, and the standard trigamma series~\cite{dlmf} gives
\[
 g'(a)=-\frac1a+(1-a)\psi'(1+a),\qquad
 \psi'(1+a)=\sum_{r=1}^{\infty}\frac1{(r+a)^2}
           <\int_a^{\infty}\frac{du}{u^2}=\frac1a.
\]
Thus $g'(a)<0$ for $0<a<1$ and consequently $g(a)>g(1)=0$.
\end{proof}

The limit $m\downarrow0$ is singular for the original process. At
$m=0$, every modulus is $1$, so a positive integer initial value never
changes and has no finite extinction time. The fixed-$m$ theorem does
not justify interchanging the limits $m\downarrow0$ and $n\to\infty$.

\section{Smoothly varying moduli}\label{rex:sec:weights}

This section proves \cref{rex:thm:weights-main} in full. The useful
information in \cref{rex:eq:smooth-weights} is the local relative increment
of the modulus; no differentiability of an interpolation is required.

\subsection{Elementary consequences of the ratio hypothesis}

\begin{lemma}[Regular variation and a summable tail]\label{rex:lem:regular}
Under \cref{rex:eq:smooth-weights}, the weights tend to infinity and are
eventually strictly increasing. For every compact interval
$[\delta,M]\subset(0,\infty)$,
\begin{equation}\label{rex:eq:weight-ratio}
 \frac{w_{\lfloor Kt\rfloor}}{w_K}\longrightarrow t^m
 \quad\text{uniformly for }\delta\leq t\leq M.
\end{equation}
For each $0<\delta\leq1$,
\begin{equation}\label{rex:eq:weight-sum}
 \frac1{K w_K}\sum_{1\leq j<\delta K}w_j
 \longrightarrow\frac{\delta^{m+1}}{m+1}.
\end{equation}
\end{lemma}

\begin{proof}
The ratio hypothesis implies
\begin{equation}\label{rex:eq:log-weight}
 \log(w_k/w_{k-1})=\frac{m+o(1)}k.
\end{equation}
In particular the ratios are eventually larger than one, and the
divergence of the harmonic series implies $w_k\to\infty$.
Summing \cref{rex:eq:log-weight} between $K$ and $\lfloor Kt\rfloor$
proves \cref{rex:eq:weight-ratio}. Uniformity follows because all relevant
indices are at least a fixed positive multiple of $K$, while the sums
of their reciprocals are uniformly bounded on $[\delta,M]$.

For the sum near zero, choose $0<\varepsilon<m$. For all sufficiently
large $\ell$,
\[
 \log(w_\ell/w_{\ell-1})\geq\frac{m-\varepsilon}{\ell}.
\]
It follows, for all sufficiently large $j\leq K$, that
\begin{equation}\label{rex:eq:power-domination}
 \frac{w_j}{w_K}
 \leq\exp\!\left(-(m-\varepsilon)
                 \sum_{\ell=j+1}^{K}\frac1\ell\right)
 \leq\left(\frac{j+1}{K+1}\right)^{m-\varepsilon}.
\end{equation}
The finitely many omitted weights have contribution tending to zero
after division by $K w_K$. On $[\eta,\delta]$ for $\eta>0$, the
locally uniform limit in \cref{rex:eq:weight-ratio} gives the usual
Riemann sum. The bound in \cref{rex:eq:power-domination} makes the
normalized sum over $j<\eta K$ uniformly small as $\eta\downarrow0$.
Thus the full limit is $\int_0^\delta t^m dt$, proving
\cref{rex:eq:weight-sum}.
\end{proof}

The proof of this lemma is included to keep the main theorem independent
of general regular-variation machinery. Its conclusion is the familiar
regular variation of index $m$, obtained here from an explicit
successive-ratio assumption.

\subsection{The same backward profile}

Define
\begin{equation}\label{rex:eq:weighted-backward}
 q_K^{(K)}=1,\qquad
 q_{k-1}^{(K)}=\ceil{\frac{w_k}{w_{k-1}}q_k^{(K)}},\qquad
 T_w(K)=q_1^{(K)}.
\end{equation}
The proof of \cref{rex:lem:duality} applies without change:
\begin{equation}\label{rex:eq:weighted-duality}
 \tau_w(n)>K\quad\Longleftrightarrow\quad n\geq T_w(K).
\end{equation}
This remains valid even if finitely many early weights decrease. The
normalization $w_1=1$ ensures that the integer initial values are exactly
the admissible quotients at stage $1$.

The exact error identity is now
\begin{equation}\label{rex:eq:weighted-error-general}
 0\leq w_{k-1}q_{k-1}-w_kq_k<w_{k-1}.
\end{equation}
It gives
\begin{align}
 w_kq_k&\leq w_K+\sum_{j=k}^{K-1}w_j,
      \label{rex:eq:weighted-global}\\
 0&\leq T_w(K)-w_kq_k<\sum_{j=1}^{k-1}w_j
      \qquad(k\geq2).\label{rex:eq:weighted-tail}
\end{align}
For $k\geq\delta K$, eventual monotonicity bounds the right-hand side
of \cref{rex:eq:weighted-global} by $(K+1)w_K$. Together with
\cref{rex:eq:weight-ratio}, this gives a uniform bound on $q_k/K$.

Set $b_k=w_k/w_{k-1}-1$. On these macroscopic intervals,
\[
 q_{k-1}-q_k=\ceil{b_kq_k},\qquad b_k=\frac{m+o(1)}k>0.
\]
Thus the increments are uniformly bounded and at least one. The
compactness, positive-drift, and ceiling-discontinuity arguments of
\cref{rex:sec:limit} apply verbatim. The limit equation is still
$y'=-C(my/t)$, and its unique selected solution is $y_m$. Therefore
\begin{equation}\label{rex:eq:weighted-profile}
 \frac{q_{\lfloor Kt\rfloor}^{(K)}}K\longrightarrow y_m(t)
 \quad\text{locally uniformly on }(0,1].
\end{equation}
The finitely many weights before monotonicity begins are outside each
fixed macroscopic interval once $K$ is large.

\subsection{Matching and inversion}

\begin{proof}[Proof of \cref{rex:thm:weights-main}]
Set $k=\lfloor\delta K\rfloor$ in \cref{rex:eq:weighted-tail}.
By \cref{rex:eq:weight-ratio,rex:eq:weighted-profile},
\[
 \frac{w_kq_k}{K w_K}\longrightarrow\delta^m y_m(\delta).
\]
By \cref{rex:eq:weight-sum}, the normalized error sum tends to
$\delta^{m+1}/(m+1)$. The same sandwich used for
\cref{rex:thm:main}, followed by $\delta\downarrow0$, gives
$T_w(K)/(K w_K)\to L_m$.

Let $J=\tau_w(n)$. Since $w_k\to\infty$, extinction is finite;
as $n\to\infty$, its index tends to infinity. Formula
\cref{rex:eq:weighted-duality} gives
\[
 T_w(J-1)\leq n<T_w(J).
\]
The ratio hypothesis implies $w_{J-1}/w_J\to1$, so both bounds are
asymptotic to $L_mJw_J$. This proves \cref{rex:eq:weights-main}.

For the forward trajectory, put $Q_k=x_k/w_k$. Survival through
$J-1$ but not $J$ gives the exact sandwich
\[
 q_k^{(J-1)}\leq Q_k<q_k^{(J)}\qquad(1\leq k\leq J-1).
\]
Equation~\eqref{rex:eq:weighted-profile} yields $Q_{\lfloor Jt\rfloor}/J
\to y_m(t)$ locally uniformly in $(0,1)$. Therefore
\[
 \frac{x_{\lfloor Jt\rfloor}}n
 =\frac{w_{\lfloor Jt\rfloor}}{w_J}\cdot
   \frac{Q_{\lfloor Jt\rfloor}/J}{n/(Jw_J)}
 \longrightarrow\frac{t^m y_m(t)}{L_m}=F_m(t).
\]
The endpoint values and the finite-mesh monotonicity argument from
\cref{rex:thm:forward-profile} again give uniform convergence on $[0,1]$.
\end{proof}

\subsection{Examples and consequences}

\begin{corollary}[Logarithmic modifiers]\label{rex:cor:log-weights}
Let $m>0$, $b\in\R$, and $A>0$. Take $w_1=1$ and
$w_k=A k^m(\log k)^b$ for $k\geq2$. Then
\begin{equation}\label{rex:eq:log-inverse}
 \tau_w(n)\asym
 \left(\frac{c_m(m+1)^b}{A}
              \frac{n}{(\log n)^b}\right)^{1/(m+1)}.
\end{equation}
The same leading result holds if the weights are changed at finitely
many stages.
\end{corollary}

\begin{proof}
Logarithmic differences give
\[
 k\log(w_k/w_{k-1})=m+\frac{b}{\log k}+o(1),
\]
so \cref{rex:eq:smooth-weights} holds. The threshold theorem gives
$A J^{m+1}(\log J)^b\asym c_m n$.
In particular $\log J\asym\log n/(m+1)$; substitution yields
\cref{rex:eq:log-inverse}. A finite change leaves the ratio hypothesis
and the asymptotic expression for $w_K$ unchanged, so the same theorem
applies.
\end{proof}

One may similarly take $w_k=A(k+\sigma)^m$ for all sufficiently large
$k$, where $k+\sigma>0$. The stopping equivalent is
$(c_m n/A)^{1/(m+1)}$. More generally, if
$w_k=k^m\ell(k)$ satisfies \cref{rex:eq:smooth-weights}, the clean inverse
statement is
\[
 \tau_w(n)^{m+1}\ell(\tau_w(n))\asym c_m n.
\]
The profile theorem remains explicit even when this inverse has no
elementary closed form.

\begin{corollary}[Regular variation of the stopping index]\label{rex:cor:inverse-rv}
Under \cref{rex:eq:smooth-weights}, for every $\lambda>0$,
\[
 \frac{\tau_w(\lfloor\lambda n\rfloor)}{\tau_w(n)}
 \longrightarrow\lambda^{1/(m+1)}.
\]
\end{corollary}

\begin{proof}
The increasing function $K\mapsto L_mK w_K$ is eventually regularly
varying with positive index $m+1$, by \cref{rex:eq:weight-ratio}.
For any constants $0<a<\lambda^{1/(m+1)}<b$, its ratios at
$\lfloor aK\rfloor$ and $\lfloor bK\rfloor$ tend respectively to
$a^{m+1}<\lambda$ and $b^{m+1}>\lambda$. Apply these inequalities
at $K=\tau_w(n)$ and use \cref{rex:eq:weights-main} and monotonicity of
the stopping index. Let $a$ and $b$ approach
$\lambda^{1/(m+1)}$.
\end{proof}

\section{Why the microscopic schedule matters}\label{rex:sec:clock}

It is tempting to regard $w_k\asym A k^m$ as all the information needed
to determine extinction. The following exact construction shows why the
local ratio condition is substantive.

\begin{proposition}[Repeating each modulus]\label{rex:prop:blocks}
Fix an integer $r\geq2$ and put
\begin{equation}\label{rex:eq:block-weights}
 w_k=\ceil{k/r}^{m}.
\end{equation}
For every integer $n>0$ and terminal stage $K\geq1$,
\begin{align}
 \tau_w(n)&=r\bigl(\tau_m(n)-1\bigr)+1,
       \label{rex:eq:block-time}\\
 T_w(K)&=T_m(\lceil K/r\rceil).
       \label{rex:eq:block-threshold}
\end{align}
Although $w_k\asym(k/r)^m$ and the weights are regularly varying
with index $m$,
\begin{equation}\label{rex:eq:block-constant}
 \frac{T_w(K)}{K w_K}\longrightarrow\frac{L_m}{r}.
\end{equation}
\end{proposition}

\begin{proof}
The operation $R_a(x)=a\lfloor x/a\rfloor$ is idempotent:
$R_a(R_a(x))=R_a(x)$. Each modulus $j^m$ occurs at the stages
$(j-1)r+1,\ldots,jr$. Only its first occurrence can change the
current value. At the initial block the modulus is $1$, so it leaves
the integer $n$ unchanged. If the unblocked process first vanishes
when modulus $J^m$ is used, the blocked process first vanishes at
stage $(J-1)r+1$. This proves \cref{rex:eq:block-time}.

Survival through stage $K$ of the blocked process is equivalent to
survival through modulus $\lceil K/r\rceil^m$ of the unblocked
process. Duality gives \cref{rex:eq:block-threshold}. Apply
\cref{rex:eq:threshold-main} and
$w_K=\lceil K/r\rceil^m$ to obtain
\[
 \frac{T_w(K)}{K w_K}
 \asym L_m\frac{\lceil K/r\rceil}{K}\longrightarrow\frac{L_m}{r}.
\]
\end{proof}

For $n=0$, both stopping times are $1$; the displayed positive-$n$
formula need not be used. In the block example the quantity
$k(w_k/w_{k-1}-1)$ is zero inside each block and approaches $rm$
at its jumps. It therefore has no limit equal to $m$.

For comparison, choose smooth weights $v_1=1$ and
$v_k=(k/r)^m$ for $k\geq2$. They have the same leading equivalent
as \cref{rex:eq:block-weights}, but satisfy the local ratio hypothesis.
Their stopping equivalent is
\[
 \tau_v(n)\asym(c_m r^m n)^{1/(m+1)},
\]
whereas \cref{rex:eq:block-time} gives
$\tau_w(n)\asym r(c_m n)^{1/(m+1)}$.
The ratio tends to $r^{1/(m+1)}$. Thus even the leading stopping
constant distinguishes the two schedules.

This example supplies a precise research boundary. A theorem for all
regularly varying moduli must include additional information about how
their growth is distributed among individual rounding steps.
